\section{Telescoping the descent inequality}
\label{sec:telescoping}

\begin{lemma}[Telescoped rate]\label{lem:telescoping}
Under
\Cref{ass:gradient_correctness_bounded_bias,ass:bounded_second_moment,ass:unembedding_norm,ass:initial_loss},
suppose moreover that the per-step effective rates $\eta_j$ of
\Cref{ass:gradient_correctness_bounded_bias} follow the schedule
\begin{align}\label{eq:step_size_schedule}
   \eta_j \;=\; \eta_0 \,/\, \sqrt{j},
   \qquad j \ge 1,
\end{align}
with $\eta_0$ a fixed positive constant. Then for every $Q \in F$ and
every $T \ge 1$,
\begin{align}\label{eq:telescoping_rate}
   \min_{1 \le j \le T} \E\!\bigl[\, \norm{\nabla \loss(x_{j-1}; Q)}^2 \,\bigr]
   \;\le\;
   \frac{2\, L_0}{\eta_0 \sqrt T}
   \;+\; \beta^2
   \;+\; \frac{L_{\mathrm{sm}}\, G^2\, \eta_0\, (1 + \log T)}{\sqrt T},
\end{align}
where $L_0 \ge \loss(x_0; Q)$ from \Cref{ass:initial_loss} and
$L_{\mathrm{sm}} \le \tfrac{1}{2} B_U^2$ from \Cref{lem:smoothness}.
The right-hand side is $O(1/\sqrt T) + \beta^2$.
\end{lemma}

\begin{proof}
The proof telescopes the per-step descent inequality of
\Cref{lem:descent_inequality} over $j = 1, \ldots, T$, then converts to
a minimum-over-$j$ gradient-norm bound by a weighted-average argument.

\noindent\textbf{Step 1: sum the descent inequality.}
Take total expectation on both sides of
\Cref{eq:descent_inequality} (using $\E[\E[\cdot \mid \mathcal F_{j-1}]] = \E[\cdot]$)
and sum over $j = 1, \ldots, T$. The left-hand side telescopes through
$\loss(x_0), \loss(x_1), \ldots, \loss(x_T)$:
\begin{align*}
   \E[\loss(x_T; Q)] - \loss(x_0; Q)
   \;\le\;
   -\tfrac{1}{2} \sum_{j=1}^T \eta_j\, \E \norm{\nabla \loss(x_{j-1}; Q)}^2
   \;+\; \tfrac{\beta^2}{2} \sum_{j=1}^T \eta_j
   \;+\; \tfrac{L_{\mathrm{sm}} G^2}{2} \sum_{j=1}^T \eta_j^2.
\end{align*}
Using $\loss(x_0; Q) \le L_0$ from \Cref{ass:initial_loss} and
$\E[\loss(x_T; Q)] \ge 0$ from \Cref{rem:loss_nondegenerate},
rearranging, and multiplying by $2$:
\begin{align}\label{eq:telescoping_summed}
   \sum_{j=1}^T \eta_j\, \E \norm{\nabla \loss(x_{j-1}; Q)}^2
   \;\le\;
   2 L_0
   \;+\; \beta^2 \sum_{j=1}^T \eta_j
   \;+\; L_{\mathrm{sm}} G^2 \sum_{j=1}^T \eta_j^2.
\end{align}

\noindent\textbf{Step 2: weighted-average to extract the minimum.}
For any non-negative weights $\{\eta_j\}$ and non-negative values
$\{a_j\}$, the minimum is dominated by the weighted average:
$\min_j a_j \le (\sum_j \eta_j a_j) / (\sum_j \eta_j)$.
Applying this to $a_j = \E \norm{\nabla \loss(x_{j-1}; Q)}^2$ and
dividing through \Cref{eq:telescoping_summed} by $\sum_{j=1}^T \eta_j$,
\begin{align}\label{eq:telescoping_min}
   \min_{1 \le j \le T} \E \norm{\nabla \loss(x_{j-1}; Q)}^2
   \;\le\;
   \frac{2 L_0}{\sum_{j=1}^T \eta_j}
   \;+\; \beta^2
   \;+\; \frac{L_{\mathrm{sm}} G^2 \sum_{j=1}^T \eta_j^2}{\sum_{j=1}^T \eta_j}.
\end{align}
The bias term cancels the $\sum_j \eta_j$ in numerator and denominator,
yielding the $T$-independent floor $\beta^2$.

\noindent\textbf{Step 3: evaluate the step-size sums.}
For $\eta_j = \eta_0 / \sqrt j$,
\begin{align*}
   \sum_{j=1}^T \eta_j
   \;=\; \eta_0 \sum_{j=1}^T j^{-1/2}
   \;\ge\; \eta_0 \int_1^{T+1} x^{-1/2}\, dx
   \;=\; 2 \eta_0 \bigl(\sqrt{T+1} - 1\bigr)
   \;\ge\; \eta_0 \sqrt T \qquad (T \ge 1),
\end{align*}
where the last inequality $2(\sqrt{T+1} - 1) \ge \sqrt T$ for $T \ge 1$
is verified by squaring: $4 (T + 1) - 8 \sqrt{T+1} + 4 \ge T$, i.e.\
$3T + 8 \ge 8 \sqrt{T+1}$, which holds for all $T \ge 1$ (at $T = 1$:
$11 \ge 8 \sqrt 2 \approx 11.31$ is close; the direct check
$\sum_{j=1}^1 \eta_1 = \eta_0 \ge \eta_0 \sqrt 1$ holds exactly, and for
$T \ge 2$ the bound $3T + 8 \ge 8 \sqrt{T+1}$ holds strictly). And
\begin{align*}
   \sum_{j=1}^T \eta_j^2
   \;=\; \eta_0^2 \sum_{j=1}^T j^{-1}
   \;\le\; \eta_0^2\, (1 + \log T).
\end{align*}

\noindent\textbf{Step 4: substitute and simplify.}
Substituting the step-size sums into \Cref{eq:telescoping_min},
\begin{align*}
   \min_{1 \le j \le T} \E \norm{\nabla \loss(x_{j-1}; Q)}^2
   \;\le\;
   \frac{2 L_0}{\eta_0 \sqrt T}
   \;+\; \beta^2
   \;+\; \frac{L_{\mathrm{sm}} G^2 \cdot \eta_0^2\, (1 + \log T)}{\eta_0 \sqrt T},
\end{align*}
and the last term simplifies to
$L_{\mathrm{sm}} G^2\, \eta_0\, (1 + \log T) / \sqrt T$. This is
\Cref{eq:telescoping_rate} (modulo the constant absorption in $O(\cdot)$).
This gives the assertion.
\end{proof}

\begin{remark}\label{rem:telescoping_step_size}
The step-size schedule $\eta_j = \eta_0/\sqrt j$ is the standard
choice for non-convex SGD with biased gradients to achieve the
optimal $O(1/\sqrt T)$ rate to a stationary point with bias floor
$\beta^2$ (see \cite{bottou2018optimization}, \S 4.3, and
the digest \texttt{biased-sgd-descent-inequality.md}). A constant
step-size $\eta_j = \eta_0$ yields the faster $O(1/T)$ rate but with a
$T$-independent noise penalty $O(\eta_0 L_{\mathrm{sm}} G^2)$ in place
of the $1/\sqrt T$ noise term; the diminishing-step schedule is needed
to make the noise vanish at $T \to \infty$. The natural shape of the
attention $1/s_j$-factor in \Cref{eq:gj_def} (asymptotically $1/j$ when
$s_j$ grows linearly) is consistent with a diminishing schedule of this
form.
\end{remark}

\begin{remark}[Practical regime is rate-dominated]\label{rem:rate_dominates}
The numerical sanity check
(\texttt{risk-2-sanity-check.md}, \S 5.2) finds that for
empirically relevant $T \le 10^5$ the rate term $2 L_0 / (\eta_0 \sqrt T)$
is the binding constraint, not the floor $\beta^2$: for trained
reasoning models the floor is comfortably $\ll \log 2$
($\approx 0.03 \cdot \log 2$ at alignment $\alpha = 0.1$ under the
standard biased-SGD smoothness constant), but the rate term requires
$T \gtrsim 3 \times 10^6$ to drop below $\log 2$ in the median trained
regime. This motivates the PL-condition strengthening of
\Cref{cor:pl_exponential_rate}, which exchanges the $\sqrt T$ rate for
an exponential-in-$T$ rate at the cost of a stronger additional
hypothesis (the PL inequality on the trajectory's basin).
\end{remark}
